% paper/sections/theorem3_identifiability.tex -- P2-04, plan/03-phase2-theory.md.
%
% Standalone section fragment, \input{} by paper/main.tex once P6-01 builds
% the skeleton. Same preamble requirements as theorem1_bound.tex /
% theorem2_lower_bound.tex. Depends on setup.tex, theorem1_bound.tex, and
% theorem2_lower_bound.tex's notation and machinery (the Fisher-information
% matrix I_k(w) from Theorem 2 Part A is reused directly, not redefined).

\section{Theorem 3: identifiability and the minimax rate}
\label{sec:theorem3}

\subsection{The rank/spacing identifiability condition}
\label{sec:theorem3-rank}

\begin{theorem}[Local identifiability: a sufficient rank condition, and a constant-rank converse]
\label{thm:identifiability}
Fix arm $k$ with true parameter $\theta_k \in \Theta$ and a design $S_{\mathrm{fit}}$ of
$m$ scales. Let $X \in \mathbb{R}^{m \times p}$ stack $J(\theta_k, s)^\top$ (Assumption
1's Jacobian) as rows, one per scale in $S_{\mathrm{fit}}$, and let
$F(\theta) = \big(g(\theta,s) - g(\theta_k,s)\big)_{s \in S_{\mathrm{fit}}}$, so
$\nabla F(\theta_k) = X$. Call $\theta_k$ \emph{locally identifiable} from noiseless
observations at $S_{\mathrm{fit}}$ if no $\theta' \ne \theta_k$ in a neighbourhood of
$\theta_k$ has $F(\theta') = 0$.
\begin{enumerate}
  \item[(a)] \emph{Sufficient.} If $\mathrm{rank}(X) = p$ (\emph{regular} local
    identifiability), then $\theta_k$ is locally identifiable.
  \item[(b)] \emph{Converse under constant rank.} If $\mathrm{rank}\,\nabla F(\theta) = r < p$
    for \emph{all} $\theta$ in a neighbourhood of $\theta_k$, then $\theta_k$ is not locally
    identifiable: the zero set of $F$ near $\theta_k$ is a $(p-r)$-dimensional submanifold
    through $\theta_k$, a continuum of observationally equivalent parameters.
  \item[(c)] \emph{No converse without a nondegeneracy hypothesis.} A rank drop at
    $\theta_k$ alone does not imply non-identifiability: for $g(\theta, s) = \theta^3$ on
    $\Theta = [-1,1]$ and $\theta_k = 0$, the parameter is uniquely identifiable from noiseless
    data, yet $J = 0$ and $\mathrm{rank}(X) = 0 < p = 1$.
\end{enumerate}
The rank condition is therefore a statement about \emph{first-order} (regular)
identifiability: full column rank certifies identifiability and gives a nonsingular Fisher
information; rank deficiency signals failure of the first-order argument, and of
identifiability itself only under the constant-rank hypothesis of~(b).
\end{theorem}

\begin{proof}
(a) By the inverse-function/injectivity theorem: $\nabla F(\theta_k) = X$ has full column
rank, so $F$ is locally injective near $\theta_k$, and $\theta_k$ is the unique zero of $F$
in a neighbourhood. (b) By the constant-rank theorem, $F$ is locally equivalent to
$(x_1, \dots, x_p) \mapsto (x_1, \dots, x_r, 0, \dots, 0)$, whose zero set is the
$(p-r)$-dimensional coordinate subspace $\{x_1 = \dots = x_r = 0\}$. (c) is the stated
counterexample: $F(\theta) = \theta^3$ vanishes only at $0$. Standard.
\end{proof}

\begin{remark}[Two different failure modes, not one]
\label{rem:two-failure-modes}
$\mathrm{rank}(X) < p$ (first-order identifiability fails; outright
non-identifiability under Theorem~\ref{thm:identifiability}(b)'s constant-rank
hypothesis) and $X$ full rank but ill-conditioned (identifiable in principle, useless in practice at finite noise) are
\emph{different} failure modes with different remedies:
\begin{itemize}
  \item \textbf{Rank deficiency} needs $m < p$ (too few scales -- $\mathrm{rank}(X) \le
    m$ always) or $m \ge p$ scales that are \emph{algebraically degenerate} for this
    particular $g$ (e.g., for \texttt{PowerLawN}/\texttt{PowerLawC}
    ($p{=}3$), two coincident scales contribute identical rows regardless of how many
    distinct scales are present elsewhere). \texttt{src/pdt/scaling/base.py}'s
    \texttt{Extrapolator.fit()} refuses $m < p+1$ outright -- \emph{stricter} than bare
    identifiability ($m \ge p$), for a reason Theorem~\ref{thm:main-bound} already
    needed: $m = p$ exactly identifies $\theta_k$ (generically) but leaves zero residual
    degrees of freedom, so $\sigma^2(s)$ and the model's own goodness of fit are
    inestimable from the fit itself.
  \item \textbf{Ill-conditioning} is a property of \emph{how} $m \ge p$ well-separated
    (in the algebraic sense above) scales are \emph{spaced}, not their count: $m$
    scales clustered close together give $X$'s rows nearly parallel (continuity of
    $J(\theta,\cdot)$ in $s$), so $X^\top X$ is nearly singular even at full formal
    rank. Section~\ref{sec:theorem3-spacing} works this out exactly for the power-law
    family.
\end{itemize}
\end{remark}

\subsection{Minimax rate}
\label{sec:theorem3-minimax}

\begin{theorem}[Local minimax rate for the in-family target $g(\theta, s^\star)$]
\label{thm:minimax}
Assume Gaussian noise with variances $\sigma^2(s)$, Assumption~\ref{asn:family}, and
Theorem~\ref{thm:identifiability}(a)'s rank condition (so $I_k(w) \succ 0$), where
$I_k(w) = \sum_s w(k,s)\, J J^\top / \sigma^2(s)$ is the information \emph{per unit of
compute} for the design's allocation $w$ ($\sum_s w(k,s)\, c(s) = 1$; Theorem~\ref{thm:lower-bound}).
Spending total compute $C$ multiplies the information by $C$, so
the local asymptotic minimax risk for estimating the \emph{in-family} quantity
$g(\theta', s^\star)$ over a shrinking neighbourhood of $\theta_k$ is, as $C \to \infty$,
\[
  \inf_{\hat\mu} \sup_{\theta' \text{ near } \theta_k} \mathbb{E}_{\theta'}\big[
    (\hat\mu - g(\theta', s^\star))^2 \big]
  \;\asymp\; \frac{1}{C}\, J(\theta_k, s^\star)^\top I_k(w)^{-1} J(\theta_k, s^\star)
  \;=\; v_k(C).
\]
This is attained, to leading order, by \emph{weighted} least squares with weights
$1/\sigma^2(s)$ (the Gauss--Markov / efficient weighting for this noise model). Unweighted
least squares under heteroscedastic noise has the larger sandwich variance
$J^\top (X^\top X)^{-1} X^\top \mathrm{diag}(\sigma^2) X (X^\top X)^{-1} J$
and matches the rate only up to a constant factor $\ge 1$, with equality when the noise is
homoscedastic.

\textbf{Two risks, not one.} The display is the risk for the \emph{in-family} target
$g(\theta', s^\star)$. If the truth is misspecified, $\mu_k(s^\star) = g(\theta_k, s^\star)
+ h_k(s^\star)$, the risk for estimating $\mu_k(s^\star)$ itself is at least
$\sigma^2_{\mathrm{extrap},k}$ plus the display, and the bias term does \emph{not} decrease with $C$
(Remark~\ref{rem:design-dependence}).
\end{theorem}

\begin{proof}[Proof sketch]
\emph{Lower bound}, via Le Cam's two-point method: for any unit direction $u \in
\mathbb{R}^p$ and perturbation size $h$, the two instances $\theta_k$ and $\theta_k +
hu$ have per-observation KL divergence at $(k,s)$ equal to $h^2 (u^\top J(\theta_k,s))^2
/ (2\sigma^2(s))$ (Theorem~\ref{thm:lower-bound} Step 3, identical calculation), so
their joint KL over the design is $h^2 u^\top I_k(w) u \, C / 2$ -- \emph{proportional to
$C$}, because $I_k(w)$ is per unit compute. Choosing $h$ so this is $O(1)$ gives
$h^2 \asymp 1 / (C\, u^\top I_k(w) u)$, and the resulting functional gap satisfies
$|g(\theta_k+hu,s^\star) - g(\theta_k,s^\star)|^2 \approx h^2 (u^\top J(\theta_k,s^\star))^2
\asymp (u^\top J)^2 / (C\, u^\top I_k u)$. Optimizing $u$ (the direction hardest to
distinguish relative to how much it moves the functional) yields
$\sup_u (u^\top J)^2 / (u^\top I_k u) = J^\top I_k^{-1} J$ and hence the stated rate
\emph{including the factor $1/C$}. \emph{This is the same change-of-measure machinery as
Theorem~\ref{thm:lower-bound} Part A, applied to estimation risk rather than selection
error}. \emph{Upper bound}: for linear $g$, weighted least squares has covariance exactly
$(C\, I_k(w))^{-1}$, hence prediction variance $J^\top I_k(w)^{-1} J / C = v_k(C)$; for
nonlinear $g$ this holds to first order (delta method), with the same non-asymptotic
caveat as Theorem~\ref{thm:main-bound}(ii). The nonlinear-$g$ statement is local and
asymptotic, not a finite-$C$ result.
\end{proof}

\subsection{Design-dependence: a worked power-law spacing example}
\label{sec:theorem3-spacing}

\begin{proposition}[Exact identifiability threshold for power-law shape parameters]
\label{prop:power-law-spacing}
For the reduced 2-parameter power law $g(A,\alpha; N) = A N^{-\alpha}$
(\texttt{PowerLawN}/\texttt{PowerLawC} minus the ceiling $E$) fit at two scales $N_1 <
N_2$, the design matrix satisfies
\[
  \det X = -A \, N_1^{-\alpha} N_2^{-\alpha} \, \log(N_2 / N_1),
\]
so $(A, \alpha)$ are identifiable from $\{N_1, N_2\}$ iff $N_1 \ne N_2$, with
conditioning governed by $|\log(N_2/N_1)|$ -- the \emph{multiplicative}, not additive,
separation of the two scales. This is why a geometric (constant-ratio) ladder, not an
arithmetic one, is the natural design for a power-law family, and why DataDecide's own
14-size ladder (4M through 1B) is close to geometric.
\end{proposition}

\begin{proof}
Direct computation: $\partial g/\partial A = N^{-\alpha}$, $\partial g/\partial\alpha =
-A N^{-\alpha}\log N$. $\det X = N_1^{-\alpha}(-AN_2^{-\alpha}\log N_2) -
N_2^{-\alpha}(-AN_1^{-\alpha}\log N_1) = -A N_1^{-\alpha}N_2^{-\alpha}(\log N_2 - \log
N_1)$, as stated. Verified numerically (not just symbolically) against
\texttt{np.linalg.det} across several $(N_1,N_2,A,\alpha)$ instances before being
trusted (see \texttt{tests/theory/test\_theorem3.py}); the two agreed to floating-point
precision in every case checked.
\end{proof}

\begin{example}[Optimal placement of a second scale is not simply "as far as possible"]
\label{ex:nonmonotone-spacing}
% NUMBER-OK: a fully worked, numerically verified illustrative example --
% the specific (A, alpha, N*, N1) values and the resulting v(N*) table are
% a computed grid search reported directly, not a claimed general result.
Fix $A=2.0$, $\alpha=0.3$, $\sigma=1$, target $N^\star = 10^9$, and $N_1 = 10^6$. A grid
search over the second scale $N_2 \in (N_1, N^\star)$, computing $v(N^\star) =
J(N^\star)^\top (X^\top X)^{-1} J(N^\star)$ exactly at each point (one replicate per
scale), gives a table that is \textbf{U-shaped, not monotone}: $v(N^\star) \approx
1.5\times 10^4$ at $N_2/N_1 = 1.01$ %NUMBER-OK: computed grid-search value, see tests/theory/test_theorem3.py::test_reproduces_the_nonmonotone_spacing_example
(nearly clustered -- Proposition~\ref{prop:power-law-spacing}'s
near-zero determinant), falling to a minimum $v(N^\star) \approx 0.52$ %NUMBER-OK: same computed grid search
near $N_2
\approx 3\times 10^7$ ($N_2/N_1 \approx 30$), then \emph{rising again} to
$v(N^\star)\approx 0.97$ as $N_2 \to N^\star$ ($N_2/N_1 = 900$). %NUMBER-OK: same computed grid search
Placing the second
scale as close as possible to the target is \emph{not} optimal for the variance at the
target, even though it minimizes that point's own extrapolation distance -- a genuine,
verified, non-obvious feature of extrapolation-design (not interpolation-design), not a
choice of illustration meant to simplify the story. This 2-point, 2-parameter case is
fully worked by hand (and checked numerically); the general $m$-point, cost-weighted,
full 3-parameter (or higher) version is the design program
$T^{\mathrm{chal}}(\nu)$ of Theorem~\ref{thm:lower-bound} Part A (a valid but not tight
characteristic time; Remark~\ref{rem:challenger-only}), with no simple closed form for
$m>2$ known to us -- P3-02's numerical solver is where the general case is actually
solved, not a hand derivation here.
\end{example}

\subsection{Numerical certificate}
\label{sec:theorem3-certificate}

\texttt{tests/theory/test\_theorem3.py} checks: (a) Proposition~\ref{prop:power-law-spacing}'s
$\det X$ formula against direct numerical computation across random
$(N_1,N_2,A,\alpha)$; (b) that estimator variance (fit via the real, production
\texttt{PowerLawN} class on many noisy replicates, not just the analytic $v_k$ formula)
blows up as design scales are clustered together and, separately, as the count of
distinct scales drops below $p+1$ (a direct \texttt{FitFailure}, not just large
variance); (c) reproduces Example~\ref{ex:nonmonotone-spacing}'s U-shaped table exactly,
confirming it is not a transcription error from the by-hand grid search above.
